% !TEX root = ../main.tex
% Abstract in English (same content as preliminares/resumen.tex).
% Goes inside \begin{abstract}[english] ... \end{abstract} in main.tex: the class adds the
% title and the key words (\tfgclaveen in datos.tex), and then the indexes.

Amyotrophic lateral sclerosis (ALS) or a stroke can leave a person unable to speak while they
still know what they want to say. In ALS, between 80 and 95\% of patients end up unable to meet
their daily communication needs through speech. The usual alternatives, such as eye tracking,
need muscles that still respond. Speech brain--computer interfaces (BCIs) take another route:
they record activity in the motor cortex while the person tries to speak and turn it into text.

In 2024, a participant with ALS and microelectrodes in the motor cortex was able to communicate
with 97.5\% word accuracy. His recordings make up the public \textit{Brain-to-Text '25}
challenge, where the best systems go below a 2\% word error rate (WER) with external language
models of hundreds of gigabytes; without them, end-to-end systems stay around 10\%. Some studies
already reuse pretrained speech models such as Whisper, but it is still unclear which part of
what they learn from audio is useful for neural signals, how they cope with new recording days
and whether they can be adapted with few resources.

This work studies these questions by adapting OWSM, an open encoder--decoder speech recognition
model, to decode intracortical neural activity using the challenge data. Its acoustic frontend
is replaced by a new one for the 512 neural features obtained every 20~ms, and everything is
trained on a single GPU without an external language model. A first block of experiments, using
all sessions, compares the pretrained encoder with the same encoder randomly initialised, the
pretrained decoder (frozen or with LoRA) with one trained from scratch, and several frontends in
terms of error and cost. The second block simulates real use: it trains models on the first
sessions in chronological order and evaluates them separately on seen and future days.

With OWSM's pretrained weights, the encoder, whether frozen or fine-tuned, does not learn to
read the signal and the WER stays at 97\%. With random weights, the same encoder reaches
20.9\%. The decoder, however, is useful: kept frozen, it brings knowledge of the language and
generalises better than one trained from scratch, and LoRA barely improves it. The best
frontend, a deep convolutional one with a frame every 80~ms, brings the WER down to 20.6\% and
trains in about three hours, 40\% less than with frames every 40~ms. In the chronological block,
the model works well on the days it has seen (WER of 20 to 27\% from 10 sessions onwards), but
on future days it stays around 80\% whether it is trained on 10 or 35 sessions. OWSM's architecture and
decoder are therefore worth reusing, and the main obstacle to real use is how the signal changes
between days.
